% !TeX program = xelatex
\documentclass[12pt,a4paper]{article}
\usepackage[margin=2cm]{geometry}
\usepackage{fontspec}
\usepackage{amsmath,amssymb,enumitem}
\setmainfont{Times New Roman}
\setlength{\parindent}{0pt}
\setlength{\parskip}{6pt}
\begin{document}
\begin{center}
{\Large\bfseries Đề thi Toán vào lớp 10 Hà Nội năm 2017 -- phần Đại số}
\end{center}

\textbf{Câu 1.} Cho hai biểu thức
\[
A=\frac{\sqrt{x}+2}{\sqrt{x}-5},\qquad
B=\frac{3}{\sqrt{x}+5}+\frac{20-2\sqrt{x}}{x-25},
\quad x\ge0,\ x\ne25.
\]
\begin{enumerate}[label=\alph*)]
  \item Tính giá trị biểu thức $A$ khi $x=9$.
  \item Chứng minh rằng $B=\dfrac{1}{\sqrt{x}-5}$.
  \item Tìm tất cả giá trị của $x$ để $A=B\cdot|x-4|$.
\end{enumerate}

\textbf{Câu 2.} Giải bài toán sau bằng cách lập phương trình hoặc hệ phương trình.

Một xe ô tô và một xe máy cùng khởi hành từ $A$ để đi đến $B$ với vận tốc của mỗi xe không đổi trên toàn bộ quãng đường $AB$ dài $120$ km. Do vận tốc xe ô tô lớn hơn vận tốc xe máy là $10$ km/h nên xe ô tô đến $B$ sớm hơn xe máy $36$ phút. Tính vận tốc của mỗi xe.

\textbf{Câu 3.}
\begin{enumerate}[label=\alph*)]
  \item Giải hệ phương trình
  \[
  \begin{cases}
  \sqrt{x}+2\sqrt{y-1}=5,\\
  4\sqrt{x}-\sqrt{y-1}=2.
  \end{cases}
  \]
  \item Trong mặt phẳng tọa độ $Oxy$, cho đường thẳng $(d):y=mx+5$.
  \begin{enumerate}[label=(\alph*)]
    \item Chứng minh đường thẳng $(d)$ luôn đi qua điểm $A(0;5)$ với mọi giá trị của $m$.
    \item Tìm tất cả giá trị của $m$ để đường thẳng $(d)$ cắt parabol $(P):y=x^2$ tại hai điểm phân biệt có hoành độ lần lượt là $x_1,x_2$ ($x_1<x_2$) sao cho $|x_1|>|x_2|$.
  \end{enumerate}
\end{enumerate}

\textbf{Câu 5.} Cho các số thực $a,b,c$ thay đổi, luôn thỏa mãn $a\ge1$, $b\ge1$, $c\ge1$ và $ab+bc+ca=9$. Tìm giá trị nhỏ nhất và giá trị lớn nhất của biểu thức
\[
P=a^2+b^2+c^2.
\]
\end{document}
